\chapter{The Circle, the Sphere, and the Cylinder}\label{ch:circle-sphere}

Area, length, and volume can be computed before defining an integral of a
function.  We start with finite polygons and prisms, then construct curved
figures by inner and outer comparisons whose gaps shrink.  Each particular
figure brings its own proof.  Neither trigonometric functions nor a general
integration theorem is needed in this chapter.

The sphere and cylinder results below are a coordinate reconstruction in
the spirit of Archimedes' exhaustion arguments, not a transcription of his
proofs.  His opening letters in \emph{On the Sphere and Cylinder} describe
the surface, volume, spherical-segment, and sector comparisons pursued
here.  The historical references are given at the end of the chapter.

\section{Finite rational geometry}
\label{thm:rational-circle-pythagorean}
The rational chart
\[
 P(t)=\left(\frac{1-t^2}{1+t^2},\frac{2t}{1+t^2}\right)
\]
satisfies $P_1(t)^2+P_2(t)^2=1$ by expansion.  For $0\le t\le1$ it
runs through the first quadrant, from $(1,0)$ to $(0,1)$.  Conversely a
point of the unit circle in this quadrant has parameter $t=y/(1+x)$.

Determinants give triangle and polygon areas.  Multiplying a polygon's
area by a rational height gives the volume of a right prism: triangulate
the base and use finite additivity.  Square-root computations give segment
lengths and hence the areas of triangles and trapezoids in rational
coordinate planes.  Computed lengths are enclosed by the rational
arithmetic of Chapter~\ref{ch:foundations}; they are not assumed to be
exact rational numbers.

Reflections and the coordinate matrix
\[
 \begin{pmatrix}c&-s\\s&c\end{pmatrix},\qquad c^2+s^2=1,
\]
preserve squared lengths and determinants by finite algebra.  Here $c,s$
are coordinates of a point, not previously defined trigonometric functions.
The broken-line inequality compares polygonal lengths.

\section{Archimedes and rational circles}
\label{thm:geometric-pi-validity}
Join finitely many points $P(t_i)$ to the origin, where
$0=t_0<\cdots<t_N=t\le1$.  For $u=t_i,v=t_{i+1}$, the inner triangle
and the outer quadrilateral cut out by the endpoint tangents have areas
\[
 I(u,v)=\frac{(v-u)(1+uv)}{(1+u^2)(1+v^2)},\qquad
 O(u,v)=\frac{v-u}{1+uv}.
\]
The formulas follow from determinants and the two tangent equations
$P(u)\mathbin{\cdot}z=P(v)\mathbin{\cdot}z=1$.
Inserting a point increases the total inner area and decreases the total
outer area: the changes are areas of finite triangles.  Also
\[
 \frac{v-u}{1+v^2}\le I(u,v)\le O(u,v)
       \le\frac{v-u}{1+u^2}.
\]
Write $\delta=\max_i(t_{i+1}-t_i)$.  The total outer--inner gap is at
most
\[
 \delta\sum_i\left(\frac1{1+t_i^2}-\frac1{1+t_{i+1}^2}\right)
 =\delta\left(1-\frac1{1+t^2}\right)\le\delta.
\]
Nested rational subdivisions therefore define a sector-area computation
$A(t)$.  Define
\[
 \pi=4A(1).
\]
A square and finer rational inscribed and circumscribed polygons give
$2<\pi<4$.  Each strict inequality is witnessed by a finite computation.
Scaling coordinates by $R$ scales every polygonal area by $R^2$; thus the
disk of radius $R$ has area computation $\pi R^2$.

\begin{figure}[htbp]\centering
\rationalCircleSubdivisionAnimation
\caption{Inner and outer rational polygons compute the circle's area.}
\label{fig:rational-parameter-line}
\end{figure}

For later use, the same inequalities applied to a subarc give
\[
 \frac{v-u}{1+v^2}\le A(v)-A(u)\le\frac{v-u}{1+u^2}
 \qquad(0\le u<v\le1).
\]
In particular $(v-u)/2\le A(v)-A(u)\le v-u$.  These are geometric
area estimates.  Their later interpretation as rectangle bounds is not
needed to construct them.

\section{Circumference and sector length}
Area and circumference are initially independent computations.  Approximate
the boundary by inscribed chords and circumscribed tangent segments.
For a unit-circle chord of length $c$, its midpoint is at distance
$d=\sqrt{1-c^2/4}$ from the origin.  Twice its radial triangle area is
$cd$, and
\[
 0\le c-cd\le c^3/4.
\]
If $c_*$ is the largest chord, the difference between the inner perimeter
and twice the inner area is at most $c_*^2$ times the inner perimeter
divided by $4$.  The perimeters have a fixed rational bound supplied by
an outer square.  The difference therefore shrinks.

For a polygon tangent to the unit circle, twice its area equals its
perimeter, triangle by triangle.  The inner and outer areas have already
been proved to agree.  The polygonal length computations consequently agree
as well, and give circumference $2\pi$.  Scaling gives circumference
$2\pi R$ at radius $R$.

The same comparison on a sector says that its arc length is twice its
area divided by its radius.  Additivity and invariance under rotation
follow by finite polygonal refinements and the determinant identities.
Computed coordinates are handled by rational enclosures and their finite
error bounds.  This proves the facts about arcs that the next chapter
will use to introduce angle parameters.

\section{Cylinders and cones by finite solid comparisons}
\label{sec:cone-exhaustion}
Extruding the circle's inner and outer polygons through a height $H$
gives inner and outer prisms.  Their volume gap is $H$ times the planar
area gap.  Thus a circular cylinder has volume computation
\[
 V_{\mathrm{cyl}}=\pi r^2H.
\]
Extruding boundary polygons instead gives lateral-area computation
$2\pi rH$.  Including both end disks gives total area
$2\pi rH+2\pi r^2$.

To compute a cone with apex at height zero, base radius $r$, and height
$H$, cut its axis at $z_k=kH/N$.  In slab $k$, the cone's radius lies
between $rk/N$ and $r(k+1)/N$.  Inscribed and circumscribed cylinder
stacks therefore have volumes
\[
 C_N^- =\frac{\pi r^2H}{N^3}\sum_{k=0}^{N-1}k^2,
 \qquad
 C_N^+ =\frac{\pi r^2H}{N^3}\sum_{k=1}^{N}k^2.
\]
These are finite solid containments, not a theorem about infinitesimal
slices.  Each circular cylinder can itself be replaced by polygonal prisms
with an arbitrarily small additional error.

The identity
$\sum_{k=1}^{N}k^2=N(N+1)(2N+1)/6$ follows by induction (or by
summing successive cubes).  Consequently
\[
 C_N^\pm=\pi r^2H\left(\frac13\mathbin{\pm}\frac1{2N}
                              +\frac1{6N^2}\right).
\]
Here the minus sign gives $C_N^-$ and the plus sign gives $C_N^+$.
Their gap is $\pi r^2H/N$.  Refining, for example with $N=2^n$, gives
nested solid bounds, and proves
\[
 \boxed{V_{\mathrm{cone}}=\frac13\pi r^2H.}
\]
All occurrences of $\pi$ in finite bounds are evaluated by its existing
rational interval computation.  The displayed expressions abbreviate
finite arithmetic on these enclosures.

\section{The hemisphere and the cone fill the cylinder in the bounds}
\label{sec:sphere-volume-exhaustion}
Fix a positive rational radius $R$.  Use height $z$ above the equatorial
plane.  The upper hemisphere has squared sectional radius $R^2-z^2$.
Compare it with the cone whose sectional radius is $z$, both inside the
cylinder of radius and height $R$.

For a rational partition $0=z_0<\cdots<z_N=R$, put
$\Delta z_i=z_{i+1}-z_i$.  Cylinder stacks bounding the hemisphere have
volumes
\[
 H_P^-=\pi\sum_i(R^2-z_{i+1}^2)\Delta z_i,
 \qquad
 H_P^+=\pi\sum_i(R^2-z_i^2)\Delta z_i.
\]
The corresponding cone stacks have volumes
\[
 C_P^-=\pi\sum_i z_i^2\Delta z_i,
 \qquad
 C_P^+=\pi\sum_i z_{i+1}^2\Delta z_i.
\]
At every finite partition there are exact identities
\[
 \boxed{H_P^-+C_P^+=\pi R^3,
        \qquad H_P^++C_P^-=\pi R^3.}
\]
Moreover their common gap is at most
\[
 \pi\max_i\Delta z_i\sum_i(z_{i+1}^2-z_i^2)
 =\pi R^2\max_i\Delta z_i.
\]
Thus the independently bounded hemisphere and cone volumes add to the
cylinder volume.  The cone calculation gives the hemisphere volume
$2\pi R^3/3$, and reflection gives
\begin{theorem}[Sphere volume]\label{thm:sphere-volume}
The sphere's solid-volume computation is
\[
 \boxed{V_R=\frac43\pi R^3.}
\]
It is two-thirds of the volume of the circumscribed cylinder of radius
$R$ and height $2R$.
\end{theorem}

For equal $N$-slab subdivisions on each hemisphere the explicit bounds are
\[
 \pi R^3\left(\frac43-\frac1N-\frac1{3N^2}\right)
 \le V_R\le
 \pi R^3\left(\frac43+\frac1N-\frac1{3N^2}\right).
\]
Their difference is $2\pi R^3/N$.  This is the entire exhaustion error,
apart from requested finite precision in the disk-area computations.
No general principle asserting equality of volumes from equality of
cross-sections has been invoked.

\section{A conical frustum's area is another polygonal computation}
\label{lem:frustum-area}
Rotate a segment with endpoint distances $r,r'$ from the axis and axial
height $h$.  Its length is $\ell=\sqrt{h^2+(r-r')^2}$.
We compute the lateral area of the resulting frustum from planar facets,
before measuring the sphere's surface.

Choose an inner polygon of the unit circle and the outer polygon given by
its endpoint tangents.  Let their areas be $a^-,a^+$, their perimeters
$p^-,p^+$, and let the inner chord lengths be $b_j$ with midpoint
distances $d_j=\sqrt{1-b_j^2/4}$.  Scale these two polygons by $r$ and
$r'$ in the two endpoint planes and join corresponding edges.

An inner trapezoid has area
\[
 T_j=\frac{r+r'}2 b_j\sqrt{h^2+d_j^2(r-r')^2}.
\]
Its slant height lies between $d_j\ell$ and $\ell$.  Thus its total
area $F^-$ satisfies
\[
 (r+r')\ell a^-
 =\frac{r+r'}2\ell\sum_j b_jd_j
 \le F^-\le\frac{r+r'}2\ell p^-.
\]
Every outer polygon edge has distance $1$ from the axis in its base
plane.  The transverse slant height of the associated outer trapezoid is
therefore exactly $\ell$.  Its total area is
\[
 F^+=\frac{r+r'}2\ell p^+=(r+r')\ell a^+.
\]
The planar comparisons $p^-\le2\pi$ and $a^-\le\pi\le a^+$ imply
\[
 F^-\le\pi(r+r')\ell\le F^+,
 \qquad F^+-F^-\le(r+r')\ell(a^+-a^-).
\]
Refinement constructs a unique lateral-area computation and proves
\[
 \boxed{\operatorname{Area}(\text{frustum})=\pi(r+r')\ell.}
\]
The case $r'=0$ is the cone's lateral area $\pi r\ell$; the case $r=r'$
is the cylinder's lateral area.  Degenerate zero-area cases are read
directly.  This proof used only planar trapezoids and the circle's
exhaustion, not an unproved formula for a surface of revolution.

\section{The finite chord identity behind the sphere's surface}
\label{lem:spherical-band-chord}
Use coordinates $(r,z)$ in a meridian half-plane, with $r\ge0$.
Take consecutive points
\[
 p=(r,z),\qquad q=(r',z'),\qquad
 r^2+z^2=(r')^2+(z')^2=R^2,\qquad z>z'.
\]
Exclude the single diameter joining both poles.  Set
\[
 \Delta z=z-z',\quad \ell=|p-q|,\quad
 m=(p+q)/2,\quad d=|m|=\sqrt{R^2-\ell^2/4}>0.
\]
The midpoint vector is perpendicular to the chord.  Its radial coordinate
is $(r+r')/2$, so elementary similar right triangles give
\[
 \frac{r+r'}2\ell=d\Delta z.
\]
Rotating the chord gives an inner frustum of area
\[
 L(p,q)=\pi(r+r')\ell=2\pi d\Delta z.
\]

For an outer comparison, take the parallel tangent to the circle at
$Rm/d$ and truncate it by the same two horizontal planes.  Its endpoint
radii exceed $r,r'$ by the same amount
\[
 k=(R-d)\frac{\ell}{\Delta z}.
\]
The slant length is still $\ell$, so the tangent frustum has area
\[
 U(p,q)=2\pi\left(d\Delta z+(R-d)\frac{\ell^2}{\Delta z}\right).
\]
These are actual finite frustum computations.  Since $d\le R$ and
$\ell\ge\Delta z$, their comparison is
\[
 \boxed{L(p,q)\le2\pi R\Delta z\le U(p,q).}
\]
The middle quantity is the independently computed area of the cylinder
band of radius $R$ between the same planes.

There is also a uniform error bound, including near the poles.  The
nonnegative radii give
\[
 (r-r')^2\le|r^2-(r')^2|\le2R\Delta z,
 \qquad \ell^2\le4R\Delta z.
\]
Using $R-d=\ell^2/(4(R+d))$ therefore yields
\[
 \boxed{U(p,q)-L(p,q)
       =\frac{\pi\ell^4}{2(R+d)\Delta z}\le2\pi\ell^2.}
\]
This explicit estimate, not an assertion about arbitrary triangulations
of a surface, will prove convergence.

\section{Sphere and zone surface areas by exhaustion}
\label{sec:sphere-area-exhaustion}
For a full meridian take, on its northern half,
\[
 (r(t),z(t))=
 R\left(\frac{2t}{1+t^2},\frac{1-t^2}{1+t^2}\right),
 \qquad t=j/N,\quad 0\le j\le N,
\]
and reflect in the equatorial plane for the southern half.  These are
rational points.  Adjacent northern chord lengths satisfy the exact identity
\[
 \ell(u,v)^2=
 \frac{4R^2(v-u)^2}{(1+u^2)(1+v^2)}\le\frac{4R^2}{N^2}.
\]
The reflected chords have the same bound.  Write $\ell_*$ for the largest
chord, so $\ell_*\le2R/N$.  Their total length is at most the meridian
semicircle length $\pi R$, already computed from polygons.

Let $L_P$ be the sum of inner frustum areas and $U_P$ the sum of the
separate tangent-frustum areas.  The latter need not join along a common
parallel; its validity is the finite comparison below, not a general claim
that containment of surfaces orders their areas.  Adding the chord
identities and the error estimates gives
\[
 L_P\le4\pi R^2\le U_P,
 \qquad U_P-L_P\le2\pi\ell_*\sum_i\ell_i
                  \le2\pi^2R\ell_*\le\frac{4\pi^2R^2}{N}.
\]
Thus $64R^2/N$ is a rational upper bound for the geometric gap.  All these
intervals overlap, independently of the chosen meridian subdivisions.
Compute the finitely many lengths, tangents, and areas to any further
prescribed rational error and take finite-prefix intersections.  This
constructs the sphere's surface-area computation, independently of its
volume, and proves
\begin{theorem}[Sphere area]\label{thm:sphere-area}
The sphere's surface-area computation is
\[
 \boxed{\mathcal S_R=4\pi R^2.}
\]
It equals the lateral area of the circumscribed cylinder, and equals
four times the area of a great-circle disk.
\end{theorem}

At the coarsest stage, the meridian consists of the two chords from the
north pole to the equator and then to the south pole.  Its bounds are
\[
 2\sqrt2\,\pi R^2\le\mathcal S_R
       \le(8-2\sqrt2)\pi R^2.
\]
Further rational subdivisions improve these bounds by the proved schedule.

The same proof works between any two fixed parallel planes with heights
$z_0<z_1$.  The chord heights telescope to $H=z_1-z_0$ instead of $2R$.
One may take rational heights and compute the endpoint radii by square
roots.  With height mesh $\delta$, the earlier bound gives
$\ell_*\le\sqrt{4R\delta}$, so refinement terminates even at a pole.
The construction proves
\[
 \boxed{\operatorname{Area}(\text{spherical zone of height }H)
                  =2\pi RH.}
\]
In particular a cap of height $h$ has area $2\pi Rh$.  The squared
straight-line distance from its pole to a point of its rim is $2Rh$.
The cap therefore has the area of the planar disk with that chord as
radius, another of Archimedes' comparisons.

Adding the two end disks to the circumscribed cylinder gives total area
$6\pi R^2$.  Thus both the sphere's volume and its surface area are
two-thirds of those of the cylinder \emph{including its bases}.  Lateral
area alone instead equals the sphere's area.  These are distinct statements.

\section{Caps and solid sectors}
\label{sec:spherical-caps-sectors}
The same finite sum identities also compute a spherical cap's volume.
Let $0\le h\le2R$ be rational and measure depth $t$ from the north pole.
The squared section radius is $q(t)=2Rt-t^2$.  For $t,u\in[0,2R]$,
\[
 |q(t)-q(u)|=|t-u|\,|2R-t-u|\le2R|t-u|.
\]
On $N$ equal depth slabs of thickness $\delta=h/N$, enclose squared
radii between $\max(0,q(j\delta)-2R\delta)$ and
$q(j\delta)+2R\delta$.  The associated cylinder stacks give volume
bounds whose gap is at most $4\pi R h\delta$.

Their reference sum, obtained using $q(j\delta)$, is exactly
\[
 \pi\delta\sum_{j=0}^{N-1}q(j\delta)
 =\pi\left[R h^2\left(1-\frac1N\right)
   -h^3\left(\frac13-\frac1{2N}+\frac1{6N^2}\right)\right].
\]
This uses only the finite sums of integers and squares.  Its difference
from $\pi(Rh^2-h^3/3)$ is bounded by
$\pi(Rh^2/N+h^3/(2N)+h^3/(6N^2))$, while either bounding stack differs
from the reference sum by at most $2\pi R h\delta$.  These bounds
construct the cap volume and prove
\[
 \boxed{V_{\mathrm{cap}}=\pi h^2\left(R-\frac h3\right).}
\]
For $h=0$ use the zero computation.  For $h=2R$ it agrees with the earlier,
independent full-sphere computation.

For $0\le h\le R$, the solid sector joining the cap to the center is
the cap together with a cone of height $R-h$ and squared base radius
$2Rh-h^2$.  The cone formula and finite additivity give
\[
 V_{\mathrm{sector}}
 =\pi h^2(R-h/3)+\frac\pi3(2Rh-h^2)(R-h)
 =\frac23\pi R^2h
 =\frac R3\operatorname{Area}(\text{cap}).
\]
For $R<h\le2R$, subtract the opposite smaller sector from the sphere;
the same formula results.  Taking the whole sphere gives
$V_R=R\mathcal S_R/3$, now as a consequence of two independent geometric
computations rather than a definition of either quantity.

\section{What precedes calculus}
Only finitely many polygons, prisms, cylinders, and planar frusta occur at
any stage.  The operations are determinants, square roots, rational sums,
and interval arithmetic.  Inner--outer gaps supply the needed number
computations.  The first cone and sphere sums have been evaluated by finite
integer identities, not by antiderivatives.  Surface area has its own
facet construction, not the rule of differentiating volume with respect
to radius.

The next chapter introduces trigonometry by using sector area as an angle
coordinate.  Only afterwards do we define a general rectangle integral
and identify these particular geometric exhaustions as examples of it.
That later organization explains and reuses the geometry; it is not a
prerequisite for the results proved here.

\begin{remark}[Historical source]
Archimedes, \emph{On the Sphere and Cylinder}, Books I and II, in T. L.
Heath's edition of \emph{The Works of Archimedes} (1897), especially the
opening letters stating the principal results.  The texts are available
as \href{https://doi.org/10.1017/CBO9780511695124.010}{Book I} and
\href{https://doi.org/10.1017/CBO9780511695124.011}{Book II}.
The rational-coordinate schedules and explicit error bounds in this
chapter are our reconstruction; they should not be attributed verbatim
to Archimedes.
\end{remark}
